% problem_id: S001-lemma-cohomological-descent-etale-fppf-modules-unbounded
% source_id: S001 (Stacks Project)
% source_file: spaces-more-cohomology.tex
% statement_locator: spaces-more-cohomology.tex lines 1488--1499
% proof_locator: spaces-more-cohomology.tex lines 1501--1622
% env: lemma
% pairing: tight (verified)
% status: complete (statement + author proof verbatim + reason + locators)
%
\begin{lemma}
\label{lemma-cohomological-descent-etale-fppf-modules-unbounded}
Let $S$ be a scheme. Let $X$ be an algebraic space over $S$.
For $K \in D_\QCoh(\mathcal{O}_X)$ the maps
$$
L\pi_X^*K \longrightarrow R\epsilon_{X, *}(La_X^*K)
\quad\text{and}\quad
K \longrightarrow Ra_{X, *}(La_X^*K)
$$
are isomorphisms. Here
$a_X : \Sh((\textit{Spaces}/X)_{fppf}) \to \Sh(X_\etale)$ is as above.
\end{lemma}

\begin{proof}
The question is \'etale local on $X$ hence we may assume $X$ is affine.
Say $X = \Spec(A)$. Then we have $D_\QCoh(\mathcal{O}_X) = D(A)$ by
Derived Categories of Spaces, Lemma
\ref{spaces-perfect-lemma-derived-quasi-coherent-small-etale-site}
and
Derived Categories of Schemes, Lemma
\ref{perfect-lemma-affine-compare-bounded}.
Hence we can choose an K-flat complex of $A$-modules
$K^\bullet$ whose corresponding complex
$\mathcal{K}^\bullet$ of quasi-coherent $\mathcal{O}_X$-modules
represents $K$.
We claim that $\mathcal{K}^\bullet$ is a K-flat complex
of $\mathcal{O}_X$-modules.

\medskip\noindent
Proof of the claim. By
Derived Categories of Schemes, Lemma
\ref{perfect-lemma-affine-K-flat}
we see that $\widetilde{K}^\bullet$ is K-flat on the scheme
$(\Spec(A), \mathcal{O}_{\Spec(A)})$.
Next, note that $\mathcal{K}^\bullet = \epsilon^*\widetilde{K}^\bullet$
where $\epsilon$ is as in Derived Categories of Spaces, Lemma
\ref{spaces-perfect-lemma-derived-quasi-coherent-small-etale-site}
whence $\mathcal{K}^\bullet$ is K-flat by
Cohomology on Sites, Lemma \ref{sites-cohomology-lemma-pullback-K-flat-points}
and the fact that the \'etale site of a scheme has enough points
(\'Etale Cohomology, Remarks \ref{etale-cohomology-remarks-enough-points}).

\medskip\noindent
By the claim we see that
$La_X^*K = a_X^*\mathcal{K}^\bullet$ and
$L\pi_X^*K = \pi_X^*\mathcal{K}^\bullet$.
Since the first part of the proof shows that the pullback
$a_X^*\mathcal{K}^n$ of the quasi-coherent module
is acyclic for $\epsilon_{X, *}$, resp.\ $a_{X, *}$, surely the proof is done
by Leray's acyclicity lemma? Actually..., no because Leray's
acyclicity lemma only applies to bounded below complexes.
However, in the next paragraph we will show the result does follow
from the bounded below case because our complex is the derived limit
of bounded below complexes of quasi-coherent modules.

\medskip\noindent
The cohomology sheaves of
$\pi_X^*\mathcal{K}^\bullet$ and $a_X^*\mathcal{K}^\bullet$
have vanishing higher cohomology
groups over affine objects of $(\textit{Spaces}/X)_\etale$ by
Lemma \ref{lemma-vanishing-adequate}.
Therefore we have
$$
L\pi_X^*K = R\lim \tau_{\geq -n}(L\pi_X^*K)
\quad\text{and}\quad
La_X^*K = R\lim \tau_{\geq -n}(La_X^*K)
$$
by Cohomology on Sites, Lemma \ref{sites-cohomology-lemma-is-limit-dimension}.

\medskip\noindent
Proof of $L\pi_X^*K = R\epsilon_{X, *}(La_X^*\mathcal{F})$.
By the above we have
$$
R\epsilon_{X, *}La_X^*K =
R\lim R\epsilon_{X, *}(\tau_{\geq -n}(La_X^*K))
$$
by Cohomology on Sites, Lemma
\ref{sites-cohomology-lemma-Rf-commutes-with-Rlim}.
Note that $\tau_{\geq -n}(La_X^*K)$ is represented by
$\tau_{\geq -n}(a_X^*\mathcal{K}^\bullet)$ which may not be the
same as $a_X^*(\tau_{\geq -n}\mathcal{K}^\bullet)$.
But clearly the systems
$$
\{\tau_{\geq -n}(a_X^*\mathcal{K}^\bullet)\}_{n \geq 1}
\quad\text{and}\quad
\{a_X^*(\tau_{\geq -n}\mathcal{K}^\bullet)\}_{n \geq 1}
$$
are isomorphic as pro-systems.
By Leray's acyclicity lemma
(Derived Categories, Lemma \ref{derived-lemma-leray-acyclicity})
and the first part of the lemma we see that
$$
R\epsilon_{X, *}(a_X^*(\tau_{\geq -n}\mathcal{K}^\bullet)) =
\pi_X^*(\tau_{\geq -n}\mathcal{K}^\bullet)
$$
Then we can use that the systems
$$
\{\tau_{\geq -n}(\pi_X^*\mathcal{K}^\bullet)\}_{n \geq 1}
\quad\text{and}\quad
\{\pi_X^*(\tau_{\geq -n}\mathcal{K}^\bullet)\}_{n \geq 1}
$$
are isomorphic as pro-systems. Finally, we put everything together as follows
\begin{align*}
R\epsilon_{X, *}La_X^*K
& =
R\epsilon_{X, *} (R\lim \tau_{\geq -n}(La_X^*K)) \\
& =
R\lim R\epsilon_{X, *}(\tau_{\geq -n}(La_X^*K)) \\
& =
R\lim R\epsilon_{X, *}(\tau_{\geq -n}(a_X^*\mathcal{K}^\bullet)) \\
& =
R\lim R\epsilon_{X, *}(a_X^*(\tau_{\geq -n}\mathcal{K}^\bullet)) \\
& =
R\lim \pi_X^*(\tau_{\geq -n}\mathcal{K}^\bullet) \\
& =
R\lim \tau_{\geq -n}(\pi_X^*\mathcal{K}^\bullet) \\
& =
R\lim \tau_{\geq -n}(L\pi_X^*K) \\
& =
L\pi_X^*K
\end{align*}
Here in equalities four and six we have used that isomorphic
pro-systems have the same $R\lim$ (small detail omitted).
You can avoid this step by using more about cohomology of the terms
of the complex $\tau_{\geq -n}a_X^*\mathcal{K}^\bullet$ proved
in Lemma \ref{lemma-vanishing-adequate} as this will prove
directly that $R\epsilon_{X, *}(\tau_{\geq -n}(a_X^*\mathcal{K}^\bullet)) =
\tau_{\geq -n}(\pi_X^*\mathcal{K}^\bullet)$.

\medskip\noindent
The equality $K = Ra_{X, *}(La_X^*\mathcal{F})$ is
proved in exactly the same way using in the final step that
$K = R\lim \tau_{\geq -n}K$ by
Derived Categories of Spaces, Lemma \ref{spaces-perfect-lemma-nice-K-injective}.
\end{proof}
